%==============================================
% VOLUME I: Zero to Bit
%==============================================
% What this volume covers:
%   The foundational question — how is information encoded in digital
%   systems at all, starting from a voltage difference?
%
%   Path: voltage → binary → numbers → characters → bytes → pointers → integrity → serialization → compression
%
% Assumes: nothing beyond basic programming familiarity
% Required by: Volume II (Silicon Horizon) and Volume III (Array Odyssey)
% CHAPTER ORDER RATIONALE (linear dependency chain):
%
%   STAGE 1  — Philosophy: what data IS before encoding
%   STAGE 2  — Physics: voltage → binary → boolean algebra → number bases
%   STAGE 3  — Integers: the first and most primitive stored value
%   STAGE 4  — Floating-point: extending integers to real numbers
%   STAGE 5  — Characters & text: non-numeric data encoding
%   STAGE 6  — Byte ordering: how multi-byte values cross system boundaries
%   STAGE 7  — Bitwise operations: the toolbox for manipulating raw bits
%   STAGE 8  — Memory alignment: how values sit in memory
%   STAGE 9  — Pointers: addresses are integers — bridge from layout to access
%   STAGE 10 — Error detection/correction: integrity before complex encodings
%   STAGE 11 — Media encodings: culmination (uses integers, floats, bitwise,
%               alignment, pointers, and error codes all at once)
%   STAGE 12 — Serialisation & compression: packaging everything for
%               storage/transmission — the closing gate to Volume II
%
%==============================================

%
\begin{volumeintro}
\lettrine[lines=3]{B}{efore we} can understand how arrays store elements, we must first understand how computers represent information itself. Every piece of data---numbers, text, images, instructions---exists as patterns of bits, and every bit begins its life as a voltage. This volume follows that single thread: from a voltage difference across a transistor, through binary encoding, to the numbers, characters, and bit-level structures that everything else in this work is built from.

\begin{quote}
\textit{``The choice of representation is often more important than the choice of algorithm.''}

\hfill--- \textsc{Donald E. Knuth}
\end{quote}
\end{volumeintro}

% ──────────────────────────────────────────────────────────────────────
% STAGE 1 — What Data Is and What It Means to Represent It
%
%   Before a single bit appears, we need to understand what data is,
%   and what it means for one physical thing to stand for another.
%   These two chapters are the philosophical foundation everything else
%   rests on.
% ──────────────────────────────────────────────────────────────────────

\chapter{The Nature of Data}
\label{ch:nature-of-data}

What data actually is before it is encoded: the distinction between data,
information, and meaning; raw measurements versus interpreted values; data
as the record of a distinction or a difference; the role of an observer or
interpreter in turning a physical trace into data; examples spanning natural
``data'' (tree rings, sediment layers, starlight) and human-made data (tally
marks, ledgers, sensor readings); why data alone is inert without a scheme
to read it; the seed of the idea that everything downstream---bits, numbers,
arrays---is data wrapped in successive layers of agreed-upon meaning.

\chapter{The Philosophy of Representation}
\label{ch:philosophy-of-representation}

What it means for one thing to stand for another; concrete versus abstract
representation and the historical shift from tally marks to abstract number
systems; Plato's allegory of the cave and the idea that representations can
be incomplete or misleading; Aristotle's more practical view of mental
representation as a useful internal model rather than a perfect copy;
medieval and modern philosophical treatments of representation (Aquinas,
Descartes, Kant, Sartre) and their relevance to encoding and interpretation;
the layered nature of representation, where symbols represent symbols all the
way down to physical reality; the abstraction hierarchy in computing, from
physical voltages to applications, and where arrays sit within it; Shannon's
information theory as the mathematical formalisation of representation,
encoding, decoding, and the bit as the fundamental unit of distinguishing
information.


% ──────────────────────────────────────────────────────────────────────
% STAGE 2 — Voltage, Binary, and Boolean Algebra
%
%   The literal starting point of computation: a transistor switching
%   between two voltage levels. From that physical fact, everything else
%   follows — binary representation, Boolean algebra, and the logical
%   operations that processors perform billions of times per second.
% ──────────────────────────────────────────────────────────────────────

\chapter{Voltage to Bits}
\label{ch:voltage-to-bits}

How transistors encode binary states; voltage levels and logic levels; noise
margins and signal integrity; why binary won over ternary and other
alternatives; the physical basis for the 0/1 abstraction that underlies
every encoding in the rest of this volume.

\chapter{Binary Representation and Boolean Algebra}
\label{ch:binary-boolean}

Two-state logic; Boolean operations (AND, OR, NOT, XOR, NAND, NOR) and their
truth tables; De Morgan's laws; the bit as the fundamental unit of
information; why Boolean algebra is the mathematics of bits; connection
to Shannon's information theory introduced in
Chapter~\ref{ch:philosophy-of-representation}.

\chapter{Number Systems and Bases}
\label{ch:number-systems}

Positional notation; decimal, binary, octal, and hexadecimal; base-conversion
algorithms; why hexadecimal is a convenient shorthand for binary; historical
development of positional systems; why different bases matter for different
purposes in systems programming.


% ──────────────────────────────────────────────────────────────────────
% STAGE 3 — Integer Representation
%
%   Numbers are the most fundamental values arrays store. Before
%   floating point, before characters, before pointers — integers.
% ──────────────────────────────────────────────────────────────────────

\chapter{Integer Representation: Unsigned}
\label{ch:unsigned-integers}

Binary positional notation for non-negative values; range limitations by
bit width; conversion algorithms; hexadecimal as a practical notation;
why unsigned integers are the natural unit of memory addressing.

\chapter{Signed Integer Representation}
\label{ch:signed-integers}

Sign-magnitude and its problems; one's complement; two's complement and why
it won (single zero, natural arithmetic, hardware simplicity); arithmetic
operations in two's complement; overflow detection; how the CPU's adder works
the same for signed and unsigned.

\chapter{Integer Overflow and Wraparound Behavior}
\label{ch:integer-overflow}

Defined wraparound versus undefined behavior (C/C++ vs.\ Rust vs.\ Java);
modular arithmetic; detecting overflow before it happens; saturating
arithmetic; real-world consequences of overflow bugs.

\chapter{Fixed-Point Representation}
\label{ch:fixed-point}

Q-format notation; scaling factors; the precision-range tradeoff; why
fixed-point dominates embedded and signal-processing systems where
floating-point hardware is absent or too slow; fixed-point arithmetic
and rounding.


% ──────────────────────────────────────────────────────────────────────
% STAGE 4 — Floating-Point Representation
%
%   Real numbers, approximation, and the cost of that approximation.
% ──────────────────────────────────────────────────────────────────────

\chapter{Floating-Point Representation: IEEE 754}
\label{ch:ieee754}

Sign, exponent, and mantissa encoding; normalised and denormalised
(subnormal) numbers; single precision (binary32) and double precision
(binary64); the implicit leading 1 bit; biased exponent; the full bit
layout of a \texttt{float} and a \texttt{double}.

\chapter{Special Floating-Point Values}
\label{ch:float-specials}

Positive and negative infinity; quiet NaN and signalling NaN; signed zero
($+0$ vs.\ $-0$); subnormal numbers and gradual underflow; how each arises
in practice and what operations produce them.

\chapter{Floating-Point Arithmetic Operations}
\label{ch:float-arithmetic}

Addition, multiplication, and division in floating-point; guard bits and
rounding; the four IEEE 754 rounding modes (round-to-nearest-even, toward
zero, toward $+\infty$, toward $-\infty$); fused multiply-add (FMA) and
why it matters for numerical accuracy.

\chapter{Floating-Point Error Analysis}
\label{ch:float-error}

Precision loss; catastrophic cancellation; machine epsilon ($\varepsilon$);
relative error versus absolute error; the Kahan compensated summation
algorithm; why naive floating-point sums of large arrays diverge from the
mathematically exact result.

\chapter{Decimal Floating-Point Representation}
\label{ch:decimal-float}

IEEE 754-2008 decimal32/decimal64/decimal128 formats; densely packed decimal
(DPD) encoding; why binary floating-point is unsuitable for financial
computation; use cases and hardware support.

\chapter{Extended Precision and Arbitrary Precision}
\label{ch:extended-precision}

The 80-bit extended-precision format (x87); quadruple precision (binary128);
arbitrary-precision arithmetic libraries (GMP, MPFR); when and why to reach
for them.


% ──────────────────────────────────────────────────────────────────────
% STAGE 5 — Character and Text Encoding
%
%   Strings are arrays of characters. Before Volume III discusses string
%   algorithms, we need to understand what a "character" actually is at
%   the byte level.
% ──────────────────────────────────────────────────────────────────────

\chapter{Character Encoding: ASCII and Extensions}
\label{ch:ascii}

7-bit ASCII; the 95 printable and 33 control characters; extended ASCII and
the proliferation of incompatible code pages (ISO 8859 family,
Windows-125x); why code pages failed for internationalisation; the need for
a universal encoding system.

\chapter{Unicode: Universal Character Encoding}
\label{ch:unicode}

Code points and the Unicode code space (U+0000 to U+10FFFF); planes and
blocks; combining characters and grapheme clusters; canonical decomposition
and normalisation forms (NFC, NFD, NFKC, NFKD); why ``number of
characters'' is not a well-defined question without specifying a counting
unit.

\chapter{Unicode Transformation Formats}
\label{ch:utf}

UTF-8: variable-width encoding, byte structure, backward compatibility with
ASCII, self-synchronisation; UTF-16: surrogate pairs and the BMP exception;
UTF-32: fixed-width simplicity at a space cost; the byte-order mark (BOM)
and when it causes problems; choosing the right encoding for a given context.

\chapter{Text Processing Complexities}
\label{ch:text-processing}

Grapheme vs.\ code point counting; case folding and locale-dependent
comparisons; text segmentation (word, sentence, line boundaries); combining
marks and their effect on string length; why ``reverse a string'' is
non-trivial in Unicode.


% ──────────────────────────────────────────────────────────────────────
% STAGE 6 — Byte Ordering and Cross-Platform Data Exchange
%
%   Every binary file format, network protocol, and serialised array
%   blob depends on a defined byte order. A short but critical stage.
% ──────────────────────────────────────────────────────────────────────

\chapter{Endianness: Byte Ordering}
\label{ch:endianness}

Big-endian vs.\ little-endian; the historical origin of the terms; network
byte order (big-endian) and why it was chosen; bi-endian architectures;
byte swapping; how endianness affects multi-byte integer and floating-point
values stored in memory or transmitted over a network.

\chapter{Cross-Platform Data Exchange}
\label{ch:cross-platform}

Portable binary formats and their endianness contracts; serialisation and
deserialisation as the mechanism for crossing byte-order boundaries; design
considerations for binary protocols that must be read on heterogeneous
systems.


% ──────────────────────────────────────────────────────────────────────
% STAGE 7 — Bitwise Operations and Bit Manipulation
%
%   The toolbox for working at the level of individual bits. Bit arrays
%   (Volume III) and many low-level techniques depend on this material.
% ──────────────────────────────────────────────────────────────────────

\chapter{Bitwise Operations Fundamentals}
\label{ch:bitwise-ops}

AND, OR, XOR, and NOT at the bit level; truth tables; bit-manipulation
primitives; using bitwise operations to implement set operations on small
universes.

\chapter{Bit Shifting and Rotation}
\label{ch:bit-shifting}

Logical shift left and right; arithmetic shift right (sign-extending);
rotate-left and rotate-right; applications: multiplication and division by
powers of 2; extracting and inserting bit fields.

\chapter{Bit Manipulation Techniques}
\label{ch:bit-techniques}

Setting, clearing, and toggling individual bits; constructing and applying
bitmasks; extracting bit fields; counting set bits (population count /
\texttt{popcount}); checking whether a value is a power of two.

\chapter{Bit Packing and Flags}
\label{ch:bit-packing}

Packing multiple Boolean flags into a single word; bitfields in structs and
their portability limitations; bit arrays as a dense Boolean store; practical
patterns from operating systems and embedded firmware.

\chapter{Advanced Bit Hacking}
\label{ch:bit-hacking}

Finding the lowest set bit; isolating and manipulating bit patterns; bit
reversal; Gray code encoding and decoding; SWAR (SIMD within a register)
techniques; why these matter for high-performance array processing.


% ──────────────────────────────────────────────────────────────────────
% STAGE 8 — Memory Alignment
%
%   Alignment determines how array elements sit in memory and how much
%   padding the compiler inserts between struct fields. Volume II
%   discusses cache-line alignment in hardware; Volume III requires
%   aligned arrays for vectorised operations.
% ──────────────────────────────────────────────────────────────────────

\chapter{Data Alignment Fundamentals}
\label{ch:alignment-fundamentals}

Natural alignment; hardware alignment requirements by type (byte, halfword,
word, doubleword, quadword); misalignment penalties on x86 and ARM;
alignment and bus transactions.

\chapter{Structure Padding and Layout}
\label{ch:struct-padding}

How the compiler inserts padding bytes between fields to satisfy each
member's alignment requirement; the rule for struct overall alignment;
reordering members to minimise padding; packing pragmas and
\texttt{\_\_attribute\_\_((packed))}; cache-line awareness in struct design.

\chapter{Alignment Optimization Techniques}
\label{ch:alignment-optimization}

Manual padding to align hot fields to cache-line boundaries;
\texttt{alignas} and \texttt{aligned\_alloc}; ensuring SIMD loads operate
on aligned addresses; the difference between a 16-byte aligned SSE load
and an unaligned load; practical patterns for high-throughput array code.


% ──────────────────────────────────────────────────────────────────────
% STAGE 9 — Pointers
%
%   A pointer is simply an aligned unsigned integer holding a memory
%   address. Having covered integers (Stage 3) and alignment (Stage 8),
%   the reader now has everything needed to understand pointers as a
%   first-class encoding, rather than as magic.
%
%   This is the bridge between the abstract model of "bytes in memory"
%   and the concrete machine model that Volume II describes in hardware.
% ──────────────────────────────────────────────────────────────────────

\chapter{Pointer Representation}
\label{ch:pointer-repr}

How a pointer is stored in memory; pointer size (32-bit vs.\ 64-bit address
spaces); the flat virtual address space model; pointer tagging — using low
bits of an aligned pointer to carry metadata; compressed pointers (JVM
\texttt{-XX:+UseCompressedOops}, V8's pointer compression).

\chapter{Pointer Arithmetic and Address Calculation}
\label{ch:pointer-arithmetic}

Computing the address of an array element from base address, element size,
and index; struct member offsets via \texttt{offsetof}; pointer subtraction
as element count; the rules that make C's address arithmetic defined and
those that make it undefined.

\chapter{Special Pointer Values}
\label{ch:special-pointers}

The null pointer and its representation; wild pointers (uninitialised) and
dangling pointers (freed or out-of-scope); pointer alignment requirements
and what a misaligned pointer means to hardware; fat pointers in Rust and
slice pointers in Go as examples of richer pointer encodings.


% ──────────────────────────────────────────────────────────────────────
% STAGE 10 — Error Detection and Correction
%
%   Data stored in arrays and transmitted across networks can be
%   corrupted. This stage covers the encoding techniques that detect and
%   correct such errors.
%
%   Placed here (before media encodings) because many container formats
%   and codecs embed CRC fields and Reed-Solomon codes directly; the
%   reader will recognise them when they appear in Stage 11.
% ──────────────────────────────────────────────────────────────────────

\chapter{Error Detection Codes}
\label{ch:error-detection}

Parity bits (even and odd parity); simple checksums; CRC algorithms and
their polynomial basis; hash-based integrity (cryptographic vs.\
non-cryptographic); tradeoffs between detection capability and overhead.

\chapter{Error Correction Codes}
\label{ch:error-correction}

The distinction between error detection and error correction; Hamming codes
and the SEC-DED (single-error-correct, double-error-detect) scheme used in
ECC RAM; Reed-Solomon codes (used in QR codes, CDs, and deep-space
communication); LDPC codes (used in Wi-Fi, 5G, and modern SSDs);
convolutional codes.

\chapter{Cyclic Redundancy Checks (CRC)}
\label{ch:crc}

CRC polynomials and their selection criteria; CRC-8, CRC-16, and CRC-32
standards; efficient implementation using precomputed lookup tables; why
hardware CRC instructions (\texttt{crc32} on x86) matter for high-speed
storage and network stacks.


% ──────────────────────────────────────────────────────────────────────
% STAGE 11 — Media Encodings
%
%   The synthesis stage: colour, images, video, and audio are all arrays
%   of encoded values. Every concept from Stages 3-10 appears here:
%
%     integers      → pixel values, sample amplitudes
%     floating-point→ HDR images, spectral analysis
%     bitwise ops   → colour packing, lossy quantisation
%     alignment     → pixel buffer layout, SIMD-friendly row strides
%     pointers      → frame buffer addressing, codec internal buffers
%     error codes   → container format checksums, codec resilience
%
%   This chapter group functions as a running integration test for the
%   reader's understanding before the final stage.
% ──────────────────────────────────────────────────────────────────────

\chapter{Color Representation Models}
\label{ch:color-models}

RGB and RGBA; HSV and HSL; CMYK; YUV and YCbCr; colour spaces and gamuts;
gamma correction and linearisation; why video and image processing pipelines
operate in different colour spaces.

\chapter{Pixel Formats and Bit Depth}
\label{ch:pixel-formats}

1-bit (monochrome); 8-bit indexed (palette); 16-bit (RGB565, RGBA4444);
24-bit RGB and 32-bit RGBA; 48-bit and 64-bit HDR formats; planar vs.\
packed layout; how pixel format choices affect array layout in memory.

\chapter{Image Compression Fundamentals}
\label{ch:image-compression}

Lossless compression (PNG, GIF, WebP lossless); lossy compression (JPEG,
WebP lossy); compression ratio; quality vs.\ file-size tradeoffs; how
compression transforms the array representation of an image.

\chapter{Video Encoding Principles}
\label{ch:video-encoding}

I-frames, P-frames, and B-frames; motion compensation and inter-frame
prediction; modern codecs (H.264/AVC, H.265/HEVC, VP9, AV1); how temporal
redundancy allows video arrays to be compressed far more than static images.

\chapter{Audio Representation: Sampling Theory}
\label{ch:audio-sampling}

The Nyquist-Shannon sampling theorem; sample rate; bit depth; quantisation
noise; aliasing and anti-aliasing filters; why 44.1\,kHz at 16 bits became
the CD standard.

\chapter{Audio Encoding Formats}
\label{ch:audio-encoding}

PCM (uncompressed); FLAC (lossless); MP3, AAC, and Opus (lossy); perceptual
coding and psychoacoustic models; why lossy compression is transparent to
human hearing at sufficient bitrates.

\chapter{Signal Processing Representations}
\label{ch:signal-processing-repr}

Time-domain representation; frequency-domain representation; the Fourier
transform as a change of array basis; spectrograms as 2-D arrays; wavelet
transforms; a preview of the FFT algorithms treated fully in Volume~III.


% ──────────────────────────────────────────────────────────────────────
% STAGE 12 — Serialisation and Compression
%
%   The closing stage: how in-memory array data is written to disk or
%   transmitted across a network (serialisation), and how it is
%   compacted to reduce storage and bandwidth costs (compression).
%
%   This is the natural handoff to Volume II (Silicon Horizon): once the
%   reader understands how data is encoded and packaged into bytes, they
%   are ready to ask "what does the hardware do with those bytes?"
% ──────────────────────────────────────────────────────────────────────

\chapter{Data Serialisation Fundamentals}
\label{ch:serialisation-fundamentals}

Converting in-memory structures and arrays to byte streams; the portability
problem (alignment, endianness, pointer size); self-describing vs.\
schema-based formats; versioning and forward/backward compatibility.

\chapter{Binary Serialisation Formats}
\label{ch:binary-serialisation}

Protocol Buffers (protobuf); FlatBuffers; Cap'n Proto; MessagePack; CBOR;
how each makes different tradeoffs between encoding speed, encoded size,
schema evolution, and random access into the serialised data.

\chapter{Text Serialisation Formats}
\label{ch:text-serialisation}

JSON, XML, YAML, and TOML; human-readable vs.\ machine-efficient tradeoffs;
parsing overhead; when text formats are appropriate and when they are not.

\chapter{Custom Binary Protocols}
\label{ch:binary-protocols}

Designing efficient custom binary formats; field alignment in wire formats;
length-prefixed vs.\ delimiter-based framing; versioning strategies; how
array payloads are typically embedded in binary protocols.

\chapter{Data Compression Theory}
\label{ch:compression-theory}

Information theory foundations: entropy and the limits of lossless
compression; the connection to Shannon's source coding theorem; lossless
vs.\ lossy bounds; why no compressor can shrink every input.

\chapter{Dictionary-Based Compression}
\label{ch:dictionary-compression}

LZ77 and its sliding-window back-references; LZ78 and its explicit
dictionary; LZW (used in GIF and early Unix \texttt{compress}); the
Lempel-Ziv family as the backbone of modern general-purpose compressors.

\chapter{Entropy Coding}
\label{ch:entropy-coding}

Huffman coding: optimal prefix-free codes, construction algorithm, canonical
Huffman codes; arithmetic coding: encoding a sequence as a single fraction;
asymmetric numeral systems (ANS) and why it replaced arithmetic coding in
zstd, brotli, and LZFSE.

\chapter{Modern Compression Algorithms}
\label{ch:modern-compression}

zlib and gzip; bzip2 (Burrows-Wheeler); LZMA and xz; Zstandard (zstd) and
its dictionary pretraining; Brotli (optimised for HTTP); compression levels
and the speed-ratio tradeoff; practical guidance for choosing a compressor.

\chapter{Specialised Compression}
\label{ch:specialised-compression}

Run-length encoding (RLE) for sparse and repetitive arrays; delta encoding
for sorted or temporally correlated sequences; columnar compression
(dictionary encoding, frame-of-reference, bit packing) used in analytical
databases; domain-specific methods for scientific data (HDF5 filters, NetCDF
compression) that appear again in Volume~III's coverage of array databases.